\section{Discussion}
\label{sec:discussion}

\subsection{Why Confirmation Resists Learning}

The two negative results of \cref{sec:inversion} and \cref{sec:learned} have a single
explanation, and it is a property of the data rather than of the methods. The confirmation
stage receives roughly two hundred positive tracks against 7000--28,000 negatives, and the
negatives are not a homogeneous class. They divide into false candidates, which are easy,
and \emph{duplicate} tracks on animals already counted, which are not: a duplicate is a
correct detection of a real animal, and \cref{tab:missed} shows that duplicates sit between
the positives and the false candidates on every feature we measured --- box size, track
length, candidate confidence alike.

A confirmer therefore receives contradictory supervision on exactly the examples that
distinguish counting from detection. That is why capacity does not help: the capacity trend
in \cref{sec:learned} degrades monotonically from 13 parameters to $6\times10^4$, and
shrinking a gradient-boosted model to depth-1 stumps recovers nearly the whole gap. It is
also why every improvement to candidate generation hurts. Growing the pool grows the
ambiguous middle of the negative class far faster than it grows the positives, so a rule
tuned to survive the new ratio must be more aggressive, and aggression costs real animals.
Extrapolating the capacity trend puts the crossover in the low thousands of positive
tracks --- an order of magnitude beyond this corpus, and a concrete target for the next
dataset rather than for a better architecture.

\subsection{Recall and Duplicate Suppression Are Bought with One Parameter}

\Cref{sec:sensitivity} makes the same mechanism visible in a single parameter and is, we
think, the most transferable observation in the paper. Deleting the confidence condition
raises held-out coverage by 23 points, more than any other change we made, and simultaneously
turns a systematic under-count into a systematic over-count, because roughly two duplicates
are admitted for each genuine animal recovered. The condition is not filtering non-animals.
It is the only cheap signal separating a first sighting from a re-sighting, and the recall it
appears to be discarding is the price of that separation.

This reframes what a practitioner should tune. The question is not ``how low can the
confidence threshold go before false positives appear'' --- the usual detection framing, under
which $c=0$ looks free, since coverage rises and the extra tracks are all on real animals. It
is ``at what point does the estimator's bias cross zero'', which is a different quantity with
a different optimum. \Cref{tab:confsens} locates that crossing at $c\approx0.45$, one full
MAE point away from the minimum-error setting. Which of the two is correct depends on the
downstream statistic: a density estimate aggregated over many transects is served by the
unbiased setting, while a per-transect report is served by the low-variance one. We report
the low-variance setting as the headline and note that the choice is a modelling decision
that the counting literature does not usually make explicit.

\subsection{Practical Guidance}

The guidance that follows is inverted from the field's default. For a counting application at
this data scale: do not select a detector by its detection metric, because the counting
criterion re-ranks seven of eleven models by two or more places (\cref{sec:choosing})
and the best detector by AP$_{50}$
counts worse than the fourth-best; do not evaluate on video the detector trained on, which
inflates coverage by 23 points (\cref{tab:leak}); do not assume a larger candidate pool is a
better one (\cref{tab:pools}); and do not expect a learned confirmation stage to beat a
hand-tuned rule until the corpus is substantially larger than 236 individuals
(\cref{tab:confirmers}).

Two positive recommendations survive the same evidence. Contrast normalization is worth more
than any architectural choice we tested --- 0.519 against 0.299 test mAP50 --- which is a
larger effect than the entire 0.365--0.510 spread across twelve architectures. And the upper
tail of the object-scale distribution deserves the attention normally reserved for the lower
one: three of our misses are animals \emph{larger} than anything in the training split
(\cref{fig:scale}), a failure mode that a small-object-focused evaluation would never
surface.

\subsection{Limitations}
\label{sec:limitations}

Several limits bound these claims. The corpus is 236 animals from four sites in one region
and one winter, one site contributes 56\% of them, per-site figures for the smallest site
(15 animals) are correspondingly noisy, and all results are single-species. Our headline
counting metric is a per-video aggregate, and the stricter identity-matched figure is 10
points lower. Four of the 13 held-out videos are the detector's validation split and so
influenced checkpoint selection, as \cref{sec:heldout} states; they carry 45 of the 83
animals. No individual animal appears in more than one video, so the re-identification limit
of \cref{sec:reid} is measured within video and cross-session identification could not be
evaluated at all. The annotated size distribution is bounded above at 96\,px, so the three
held-out animals larger than 100\,px lie outside it (\cref{fig:scale}), and the natural
remedy is unavailable in the symmetric form the standard augmentation provides. Finally, the
sensitivity analysis of \cref{sec:sensitivity} varies one parameter of one rule on one
operating pool; it establishes that the confidence condition is load-bearing here, not that
0.65 transfers to another sensor or another species.
